\cleardoublepage
\pagenumbering{arabic}
\pagestyle{plain}

\chapter*{Introducción}
\addcontentsline{toc}{chapter}{\textnormal{Introducción}}
\label{chap:introduccion}

\setstretch{1.5}
\fontsize{12}{14}\selectfont
\vspace{\baselineskip}



No cabe duda de que, en el ámbito de la ingeniería lingüística, la terminología es una materia de gran importancia, tanto en su vertiente descriptiva como aplicada (\cite{cabre2004termintegral}). Bajo esta premisa, la disciplina se ha consolidado dentro de las ciencias del lenguaje contemporáneas como una materia transversal, cuya relevancia no se limita a la descripción de los lenguajes especializados, sino que ha avanzado a pasos agigantados hacia el desarrollo de herramientas digitales y computacionales para el tratamiento, gestión y sistematización del léxico de especialidad. A medida que los lenguajes técnicos se diversifican y evolucionan, se vuelve necesario no solo registrar las unidades que los constituyen, sino también comprender y dar cuenta de aquellas propiedades que hacen que una palabra sea término. En este contexto, la intersección entre teoría lingüística, lingüística de corpus y tecnologías del lenguaje ofrece un terreno fértil para explorar nuevas formas de detección, formalización y modelización de unidades terminológicas.

En esta tesis doctoral proponemos la extracción automática de unidades terminológicas mediante la combinación de enfoques lingüísticos, estadísticos y computacionales, con el objetivo de avanzar hacia una modelización teórica y computacional de la noción de 'unidad terminológica' o 'término'. Para ello, trabajamos sobre un corpus especializado en lingüística de corpus, por tratarse de un subdominio que presenta una elevada densidad terminológica, una rápida evolución conceptual y una notable interdisciplinariedad. Estas características lo convierten en un entorno idóneo tanto para la observación empírica de unidades terminológicas en contexto como para la validación de estrategias de extracción automática y abstracción formal. Además, la elección de este dominio facilita la integración de métodos computacionales con principios teóricos procedentes de la propia disciplina lingüística, lo que permite articular un enfoque genuinamente interfacial entre la teoría terminológica y las tecnologías del lenguaje.

El dinamismo y dispersión que caracterizan a la terminología lingüística han sido siempre un problema constante y de larga tradición para los lingüistas (Perea, 2000), parejo a su interrelación con distintos campos del saber y su amplia variedad de líneas de investigación, especialmente en los últimos años. La consolidación de la lingüística como ciencia ha desembocado en una pluralidad de términos al que ya se apunta desde finales del siglo pasado:

\begin{quote}
    El desarrollo de la terminología lingüística ha ido de la mano del desarrollo del pensamiento lingüístico; a medida que se crean y se especifican ideas, hipótesis y teorías, la terminología se pluraliza, aumenta; la pluralidad terminológica refleja la diversidad teórica o de teorías dentro de las concepciones lingüísticas en un doble sentido, por un lado, en el de la relación entre teorías y terminologías individuales o particulares y, por otro lado, en el de la relación entre uso de términos diferentes para hacer referencia a las mismas realidades y nociones. (\cite{caballero1994reflexiones})
\end{quote}

Desde su establecimiento como ciencia a principios del siglo XX, la lingüística ha contado con un gran número de diccionarios y glosarios de consulta que tratan de consignar su extensa terminología (\cite{laborda2017diccionarios}). Estas obras lexicográficas, sin embargo, presentan una serie de cuestiones problemáticas o que no llegan a satisfacer nuestros intereses desde un punto de vista terminológico (i.e. los métodos de adición de unidades terminológicas al diccionario, sin tener en cuenta enfoques con evidencias de corpus; la concepción de la lingüística en un sentido amplio, sin tener en cuenta ramas de la lingüística aplicada o léxicos especializados centrados en dichas ramas). La ausencia de algunos términos en la compilación de diccionarios es una cuestión que ya se señala como una de las razones por las que el trabajo terminográfico se base actualmente en corpus: “[t]he usual consultation of sources [...] do not always compile all terms that we need” (\cite{oliver2007}).

El planteamiento central de esta tesis doctoral sostiene que la combinación de enfoques lingüísticos, estadísticos y computacionales —incluyendo métodos de aprendizaje automático y aprendizaje profundo— permite mejorar significativamente la extracción automática de términos en ámbito especializados, tanto en recuperación como en precisión, frente a métodos tradicionales. Previo al desarrollo de la investigación, se plantean las siguientes hipótesis:

\begin{enumerate}
    \item  La integración sinérgica de enfoques lingüísticos y estadísticos en los sistemas automáticos de extracción de términos permite una mejora sustancial tanto en la recuperación como en la precisión de los términos identificados en corpus técnicos. A diferencia de los métodos más rígidos —ya sean puramente lingüísticos, centrados en patrones morfosintácticos y estructuras gramaticales, o puramente estadísticos, enfocados en medidas de asociación o frecuencia—, los enfoques híbridos permiten captar tanto las regularidades formales como las distribuciones empíricas de los términos en contexto. Esta combinación metodológica resulta especialmente eficaz en dominios técnicos y científicos donde la variación terminológica y la alta densidad informativa dificultan la extracción mediante técnicas unidimensionales. Asimismo, se argumenta que este enfoque integrador se alinea con las propiedades  de los términos, cuya identificación computacional requiere modelos que capturen simultáneamente su forma, uso y función discursiva.
    \item La caracterización automática de unidades terminológicas requiere incorporar información más allá de propiedades léxicas y morfosintácticas aisladas. Se sostiene que la integración de información relacional —derivada del comportamiento de la unidad en su entorno discursivo, sus asociaciones recurrentes y su funcionalidad dentro del texto especializado— mejora significativamente la detección y clasificación de términos frente a enfoques basados únicamente en patrones formales o frecuencias locales. Esta información relacional es capturada más eficazmente por modelos que operan a nivel de secuencia y contexto que por sistemas que procesan tokens de forma aislada.
    \item  Los modelos computacionales basados en aprendizaje automático y aprendizaje profundo, particularmente aquellos que incorporan representaciones distribucionales de palabras (como \textit{word embeddings}, \textit{contextual embeddings} o modelos \textit{Transformer}), superan en eficacia a los métodos tradicionales de extracción de términos, tanto en tareas de identificación como en clasificación y normalización terminológica. Esta superioridad se manifiesta especialmente cuando los modelos se entrenan sobre datos lingüísticamente enriquecidos, es decir, textos anotados con etiquetas morfosintácticas, funciones discursivas, límites terminológicos y otras propiedades lingüísticas relevantes. Se argumenta que estas técnicas permiten modelar computacionalmente la noción de término como una construcción situada, sensible al contexto y potencialmente polimorfa, lo cual no puede lograrse mediante métodos puramente formales o estadísticos. Además, el aprendizaje automático facilita la incorporación dinámica de nuevas evidencias lingüísticas durante el entrenamiento, lo que resulta crucial en dominios altamente especializados y en evolución.
    \item  Los resultados obtenidos mediante técnicas avanzadas de extracción automática de términos van más allá de su valor instrumental o auxiliar —al facilitar la creación de glosarios, índices o bases de datos terminológicas—, sino que también permiten generar inferencias de orden teórico sobre la naturaleza, el comportamiento y la organización de los términos en uso real. En particular, se argumenta que la modelización computacional de patrones terminológicos basada en datos empíricos posibilita la formulación de hipótesis lingüísticas más sólidas sobre los mecanismos de formación, variación, especialización y difusión de los términos en contextos técnicos y científicos. Esta retroalimentación entre el análisis computacional y la reflexión teórica permite el desarrollo de una terminología computacional con fundamento lingüístico, en la que los modelos automáticos no solo operan como herramientas de detección, sino como instrumentos epistémicos para la comprensión profunda del fenómeno terminológico. De este modo, se propicia una integración entre teoría y práctica que fortalece el carácter científico y aplicable de la disciplina terminológica.
    \item  El grado de representatividad terminológica de los corpus empleados en la extracción automática de términos constituye un factor determinante en la calidad y fiabilidad de los resultados obtenidos, hecho fundamentado en la premisa de que la distribución y densidad de los términos en un corpus está directamente condicionada por su capacidad de reflejar de forma equilibrada y exhaustiva los distintos subdominios, géneros discursivos y niveles de especialización del ámbito temático en cuestión. Se sostiene que corpus deficientemente representativos —ya sea por sesgo temático, escasez de fuentes, baja diversidad textual o insuficiente cobertura léxica— generan un sesgo en la identificación automática de términos, afectando tanto a la variedad como a la pertinencia de las unidades extraídas. En consecuencia, se argumenta la necesidad de aplicar métricas compuestas de representatividad terminológica (como el índice de cobertura terminológica o el β de Heaps) que permitan evaluar y garantizar la calidad del corpus antes de su explotación computacional.
    \item  La incorporación de mecanismos de explicabilidad en los modelos computacionales empleados para la extracción de términos es esencial para garantizar la interpretabilidad lingüística de los resultados y, por tanto, su integración en una terminología computacional con fundamento teórico. A diferencia de los modelos puramente “caja negra”, los enfoques explicables permiten rastrear los factores que conducen a la identificación de una unidad como término, ya sea a través de rasgos morfosintácticos, medidas estadísticas o patrones contextuales aprendidos. Esta capacidad explicativa no solo favorece la validación manual y la confianza en el sistema, sino que también facilita la retroalimentación entre los resultados computacionales y la teoría terminológica, permitiendo afinar conceptos como variación, saliencia o grado de terminologización. Por ello, se defiende que los sistemas de extracción deben diseñarse con criterios de transparencia, trazabilidad y replicabilidad.
\end{enumerate}

Para la comprobación de las hipótesis expuestas, se proponen los siguientes objetivos generales y específicos, estructurados en torno a los tres ejes fundamentales que articulan esta investigación: la formalización teórico-computacional del concepto de término, el desarrollo y evaluación de métodos automáticos para su detección, y la sistematización de los resultados en un recurso terminológico estructurado. Los objetivos propuestos requieren el diseño y la aplicación de un flujo metodológico secuencial, tal y como se describe en la Figura \ref{fig:flujo_simple_colores}.

Cada objetivo general se acompaña de un conjunto de objetivos específicos que orientan las etapas del trabajo, desde la construcción del corpus hasta la validación de los modelos y la explotación de los resultados, con el fin de garantizar la coherencia entre el planteamiento teórico, la metodología empleada y los resultados obtenidos.

\begin{itemize}
    \item \textbf{O. G. 1.}: Formalizar computacionalmente la noción de término especializado a partir de evidencias empíricas observables en corpus representativos y lingüísticamente anotados.
    \begin{itemize}
        \item \textit{O. E. 1. 1}: Diseñar y construir un corpus especializado representativo del dominio de la lingüística, segmentado y enriquecido con anotaciones morfosintácticas, discursivas y terminológicas, y evaluar su representatividad terminológica y adecuación a la tarea terminográfica propuesta.
        \item \textit{O. E. 1. 2.}: Identificar, sistematizar y operacionalizar un conjunto de rasgos lingüísticos, sintáctico-semánticos, contextuales y pragmáticos relevantes para la caracterización de las unidades terminológicas.
        \item \textit{O. E. 1. 3.}: Elaborar un modelo teórico-computacional que integre los rasgos seleccionados en una abstracción formal del término, compatible con tareas de extracción y análisis automatizado.
    \end{itemize}
    \item \textbf{O. G. 2.}: Desarrollar, entrenar y evaluar sistemas automáticos de detección de términos especializados, basados en metodologías lingüísticas, estadísticas, híbridas y de aprendizaje supervisado y profundo.
    \begin{itemize}
        \item \textit{O. E. 2. 1.}: Implementar una batería de métodos de extracción terminológica que incluya enfoques basados en reglas, medidas estadísticas clásicas, métodos híbridos y modelos de aprendizaje automático y profundo.
        \item \textit{O. E. 2. 2}.: Diseñar un protocolo de validación que contemple la exclusión sistemática de términos del conjunto de entrenamiento, con el fin de medir la capacidad de generalización y la detección de términos no vistos.
        \item \textit{O. E. 2. 3.}: Evaluar el rendimiento de los modelos mediante métricas estándar en el ámbito del procesamiento del lenguaje natural, complementadas con indicadores de coherencia terminológica y explicabilidad del modelo.
        \end{itemize}
    \item \textbf{O. G. 3.}: Construir un glosario digital de términos especializados contextualizados, estructurado a partir de los resultados del sistema de extracción y orientado a aplicaciones en investigación terminológica.
    \begin{itemize}
        \item \textit{O.E . 3. 1.}: Diseñar un modelo de entrada terminológica que integre definición, variantes, ejemplos en contexto, relaciones semánticas y metadatos lingüísticos derivados del análisis computacional.
        \item \textit{O. E. 3. 2.}: Generar el glosario a partir de las unidades terminológicas extraídas y validadas, incorporando mecanismos de trazabilidad hacia el corpus fuente y el modelo utilizado.
        \item \textit{O. E. 3. 3.}:  Vincular el glosario con los resultados del modelo y con los rasgos lingüístico-computacionales que motivaron su identificación, garantizando así la trazabilidad terminológica, la replicabilidad metodológica y la posibilidad de ampliar el recurso.
    \end{itemize}
\end{itemize}

\begin{figure}[H]
\centering
\begin{tikzpicture}[node distance=4.2cm, every node/.style={font=\small}]

% Etapas principales con colores diferentes
\node (modelo) [rounded corners=3pt, minimum width=4cm, minimum height=1.2cm, align=center, fill=violet!20] {Modelo teórico-\\computacional};
\node (extraccion) [right of=modelo, rounded corners=3pt, minimum width=4cm, minimum height=1.2cm, align=center, fill=teal!20] {Extracción automática\\de términos};
\node (glosario) [right of=extraccion, rounded corners=3pt, minimum width=4cm, minimum height=1.2cm, align=center, fill=orange!20] {Glosario computacional\\contextualizado};

% Flechas principales
\draw[flecha] (modelo) -- (extraccion);
\draw[flecha] (extraccion) -- (glosario);

\end{tikzpicture}
\caption{Flujo metodológico secuencial de la investigación}
\label{fig:flujo_simple_colores}
\end{figure}

De este modo, la tesis articula un recorrido metodológico que no solo avanza linealmente desde la teoría hacia la aplicación, sino que incorpora mecanismos de retroalimentación entre los resultados obtenidos y la formalización conceptual, permitiendo una revisión progresiva del modelo a partir de los datos empíricos. Este enfoque integrador pone de manifiesto la estrecha relación entre teoría y datos, al defender que, si bien una base teórica sólida es condición necesaria para la eficacia de los modelos computacionales, los datos empíricos también desempeñan un papel activo en la reformulación de hipótesis y en la comprensión profunda del lenguaje y las lenguas. En última instancia, esta interacción dinámica responde a uno de los fines centrales de la investigación lingüística: observar, describir y explicar el comportamiento real del lenguaje y las lenguas a partir de su uso, combinando intuición teórica y evidencia cuantitativa. Esta articulación resulta aún más pertinente en el desarrollo de modelos computacionales, donde el objetivo no es únicamente detectar regularidades, sino también formalizar, explicitar y modelar de forma sistemática determinados fenómenos lingüísticos a partir de datos observables.

La tesis se estructura en seis capítulos, además la introducción y las conclusiones finales. El Capítulo \ref{chap:1_terminologia} está dedicado al estudio de la terminología desde una orientación lingüística, con especial atención a la noción de término o unidad terminológica, su variabilidad, su comportamiento discursivo en textos especializados y su representación en productos lexicográficos y terminográficos. El Capítulo \ref{chap:2_tecnologias} contextualiza la investigación en el marco de las tecnologías del lenguaje, abordando tres áreas clave —lingüística computacional, lingüística de corpus e ingeniería lingüística— cuya intersección resulta esencial para el desarrollo de modelos automáticos de extracción terminológica. En el caso de la lingüística de corpus, este capítulo permite, además, delimitar conceptualmente el dominio sobre el que se construye el corpus de trabajo. El Capítulo \ref{chap:3_extraccion} expone la extracción automática de términos, presentando sus fundamentos, los principales enfoques existentes en la literatura y su relación con los desafíos terminológicos descritos previamente. La metodología de trabajo se detalla en el Capítulo \ref{chap:4_metodologia}, incluyendo la construcción y anotación del corpus, la definición de rasgos, la implementación de modelos y los criterios de evaluación empleados. En el Capítulo \ref{chap:5_resultados} se presentan y analizan los resultados obtenidos, tanto en términos de rendimiento de los sistemas como de coherencia terminológica de las unidades detectadas. Finalmente, se recogen las conclusiones generales, se discuten las implicaciones teóricas y metodológicas derivadas del presente estudio y se proponen líneas futuras de investigación.